\documentclass[12pt]{article}

% Layout
\usepackage[a4paper,includeheadfoot,margin=2.54cm]{geometry}

% Spracheinstellungen, alle Sprachen laden, letzte ist aktiv
\usepackage[english,ngerman]{babel}

% hilfreiche Pakete der AMS (American Mathematical Society) laden
\usepackage{amsmath}
\usepackage{amssymb}

% Sätze, Lemmata, ..
\usepackage{amsthm}
\newtheorem{theorem}{Satz}[section]
\newtheorem{lemma}[theorem]{Lemma}
\theoremstyle{definition}
\newtheorem{definition}[theorem]{Definition}

% Formelnummerierung
\numberwithin{equation}{section}

% moderne Literaturverwaltung mittels Biber, erstzt BibLaTeX,
% kann UTF-8
\usepackage{biblatex}
\addbibresource{Bachelorarbeit_Vorname_Nachname.bib}

% Einbinden von Grafik, neue Version
\usepackage{graphicx}

% Unterabbildungen
\usepackage{subcaption}

% TikZ ist kein Zeichenprogramm (doch)
\usepackage{tikz}

% listings bindet Code ein
\usepackage{listings}
\definecolor{hellgrau}{rgb}{0.90,0.90,0.90}
\definecolor{commentcol}{rgb}{0.0823,.4902,0.0}
\lstset{language=Python,
        basicstyle={\footnotesize\ttfamily},
        keywordstyle={\sffamily\bfseries},
        tabsize=2,
        numbers=left,
        numberstyle=\tt,
        stepnumber=1,
        numbersep=7pt,
        breaklines=true,
        frame=single,
        frameround=ffff,
        commentstyle=\color{commentcol},
        backgroundcolor=\color{hellgrau},
        literate=
  {á}{{\'a}}1 {é}{{\'e}}1 {í}{{\'i}}1 {ó}{{\'o}}1 {ú}{{\'u}}1
  {Á}{{\'A}}1 {É}{{\'E}}1 {Í}{{\'I}}1 {Ó}{{\'O}}1 {Ú}{{\'U}}1
  {à}{{\`a}}1 {è}{{\`e}}1 {ì}{{\`i}}1 {ò}{{\`o}}1 {ù}{{\`u}}1
  {À}{{\`A}}1 {È}{{\'E}}1 {Ì}{{\`I}}1 {Ò}{{\`O}}1 {Ù}{{\`U}}1
  {ä}{{\"a}}1 {ë}{{\"e}}1 {ï}{{\"i}}1 {ö}{{\"o}}1 {ü}{{\"u}}1
  {Ä}{{\"A}}1 {Ë}{{\"E}}1 {Ï}{{\"I}}1 {Ö}{{\"O}}1 {Ü}{{\"U}}1
  {â}{{\^a}}1 {ê}{{\^e}}1 {î}{{\^i}}1 {ô}{{\^o}}1 {û}{{\^u}}1
  {Â}{{\^A}}1 {Ê}{{\^E}}1 {Î}{{\^I}}1 {Ô}{{\^O}}1 {Û}{{\^U}}1
  {Ã}{{\~A}}1 {ã}{{\~a}}1 {Õ}{{\~O}}1 {õ}{{\~o}}1
  {œ}{{\oe}}1 {Œ}{{\OE}}1 {æ}{{\ae}}1 {Æ}{{\AE}}1 {ß}{{\ss}}1
  {ű}{{\H{u}}}1 {Ű}{{\H{U}}}1 {ő}{{\H{o}}}1 {Ő}{{\H{O}}}1
  {ç}{{\c c}}1 {Ç}{{\c C}}1 {ø}{{\o}}1 {å}{{\r a}}1 {Å}{{\r A}}1
  {€}{{\euro}}1 {£}{{\pounds}}1 {«}{{\guillemotleft}}1
  {»}{{\guillemotright}}1 {ñ}{{\~n}}1 {Ñ}{{\~N}}1 {¿}{{?`}}1
  {·}{{$\cdot$}}1
}

% Hyperlinks in Texten
\usepackage{hyperref}
\hypersetup{%
  pdftitle     = {Modeling of stochastic optimization algorithms with stochastic differential equations},
  pdfsubject   = {Masterarbeit von Jonathan Hellwig},
  pdfkeywords  = {Masterarbeit, neuronale Netze},
  pdfauthor    = {\textcopyright\ Jonathan Hellwig 2022},
  linkcolor    = red,     % links to same page
  urlcolor     = blue,     % links to URLs
  citecolor    = green!50!black,     % links to citations
  breaklinks   = true,       % links may (line) break
  colorlinks   = true,
  citebordercolor=0 0 0,  % color for \cite
  filebordercolor=0 0 0,
  linkbordercolor=0 0 0,
  menubordercolor=0 0 0,
  urlbordercolor=0 0 0,
  pdfhighlight=/P,   % moeglich /I, /P, ...
  pdfborder=0 0 0,   % keine Box um die Links!
}
% nützliche Kurzkommandos für natürliche, ..., reelle, .. Zahlen
\newcommand{\N}{\mathbb{N}}
\newcommand{\Z}{\mathbb{Z}}
\newcommand{\Q}{\mathbb{Q}}
\newcommand{\R}{\mathbb{R}}
\newcommand{\C}{\mathbb{C}}
\newcommand{\F}{\mathbb{F}}
\newcommand{\K}{\mathbb{K}}

% ein paar dick gedruckte Buchstaben
\newcommand{\bfA}{\mathbf{A}}
\newcommand{\bfB}{\mathbf{B}}
\newcommand{\bfa}{\mathbf{a}}
\newcommand{\bfb}{\mathbf{b}}

% Beginn des Dokumentes
\begin{document}

% Titelseite
\thispagestyle{empty}

\begin{center}
  \includegraphics[height=2.3cm]{pics/MATH_de}
  \hfill
  \includegraphics[height=2.3cm]{pics/TUHH_de}
\end{center}

\vspace*{5em}

\begin{center}
  {\Huge
    \textsc{Titel der Bachelorarbeit}\\[2em]
  }
  {\LARGE
    Bachelorarbeit
  }

  \vspace*{2em}

  {\Large
    von\\
    Vorname Nachname\\
    aus Ort\\
    Matrikelnummer: 4711\\
    Studiengang: Technomathematik\\
  }
\end{center}

\vfill
\begin{center}
  \today
\end{center}
\vfill

\begin{tabbing}
  % längste Zeile zuerst duplizieren, mit kill löschen
  Erstprüferin: \= Prof. Dr. Sabine Le Borne\kill
  Erstprüferin: \> Prof. Dr. Sabine Le Borne\\
  Zweitprüfer:  \> Dr. Jens-Peter M. Zemke\\
  Betreuer:     \> Dr. Jens-Peter M. Zemke\\
\end{tabbing}
\newpage
% this page intentionally left blank
\thispagestyle{empty}
\mbox{}
\newpage

\section*{Eidestattliche Erklärung}

Hiermit versichere ich an Eides statt, dass ich die vorliegende
Bachelorarbeit mit dem Titel
\begin{quote}
  „Titel der Bachelorarbeit“  
\end{quote}
selbständig und ohne unzulässige fremde Hilfe verfasst habe. Ich habe
keine anderen als die angegebenen Quellen und Hilfsmittel benutzt,
sowie wörtliche und sinngemäße Zitate kenntlich gemacht. Die Arbeit
hat in gleicher oder ähnlicher Form noch keiner Prüfungsbehörde
vorgelegen. Ich versichere, dass die eingereichte schriftliche Fassung
der auf dem beigefügten Medium gespeicherten Fassung entspricht.

\vspace*{3em}

\begin{tabbing}
  \rule{.4\textwidth}{1pt} \hspace*{.2\textwidth}
  \= \rule{.4\textwidth}{1pt} \\
  Ort, Datum \> Unterschrift
\end{tabbing}

\newpage
\mbox{}
\newpage
% TOC - Table of Contents
\tableofcontents
\newpage
\listoffigures
\newpage

\section{Motivation}
\label{sec:Motivation}

\section{Background on Optimization}
\label{sec:Optimization}

\subsection{Formal definition, Existence, Uniqueness}
\subsection{Iterative methods}
\subsubsection{Gradient Descent Method}
\subsubsection{Newton's method}
\subsection{Stochastic Optimization}
\label{sec:StochasticOptimization}
\subsubsection{Stochastic Gradient Descent}
% Convergence results
% Variants
% Gradient flow

\subsubsection{State of the art algorithms}
\label{sec:StateOfTheArtAlgorithms}

\section{Background SDE theory}
\label{sec:BackgroundSDETheory}
\subsection{Ito Integral}
\label{subsec:ItoIntegral}
\subsection{SDE Definition}
\label{subsec:SDEDefinition}
\subsection{SDE Existence Uniqueness}
\label{subsec:SDEExistenceUniqueness}
\subsection{SDE Numerical methods}

\section{Modeling stochastic model with SDE}
\subsection{SGD}
\subsection{Generalization}
\section{Solving SDE Model}
\section{}
\printbibliography

\end{document}
